Large language models (LLMs) are increasingly leveraged to empower autonomous agents to simulate human beings in various fields of behavioral research.
However, evaluating their capacity to navigate complex social interactions remains a challenge.
Previous studies face limitations due to insufficient scenario diversity, complexity, and a single-perspective focus.
To this end, we introduce \textbf{AgentSense}: Benchmarking Social Intelligence of Language Agents through Interactive Scenarios.
Drawing on Dramaturgical Theory, AgentSense employs a bottom-up approach to create 1,225 diverse social scenarios constructed from extensive scripts.
We evaluate LLM-driven agents through multi-turn interactions, emphasizing both goal completion and implicit reasoning.
We analyze goals using ERG theory and conduct comprehensive experiments.
Our findings highlight that LLMs struggle with goals in complex social scenarios, especially high-level growth needs, and even GPT-4o requires improvement in private information reasoning. Code and data are available at ~\url{https://github.com/ljcleo/agent_sense}.